\section{$k$-positive maps, Schmidt rank, and Choi compression}

This section develops the $k$-positive-map criteria of
\cite[Chapter~3]{Wolf2012Quantum}. Schmidt-rank bounds convert positivity of
ampliations into quadratic-form conditions on the Choi matrix, while Ky
Fan's maximum principle controls the extremal overlap with vectors of bounded
Schmidt rank. The right-factor compression criteria give rectangular and
rank-$k$ projection forms of the same condition.

\subsection{Two-positive maps and the generalized Schwarz inequality}

\begin{definition}[$k$-positive map]\label{def:k_positive_map}
    \lean{IsNPositiveMap}
    \leanok
    A linear map
    $E : \MN{D_{\rm in}} \to \MN{D_{\rm out}}$ is $k$-positive if
    $E \otimes \id_{k}$ is positive from
    $\MN{D_{\rm in}} \otimes \MN{k}$ to
    $\MN{D_{\rm out}} \otimes \MN{k}$.
\end{definition}

\begin{definition}[Two-positive map]
    \label{def:2positive_map}
    \lean{Is2PositiveMap}
    \leanok
    A linear map is \emph{two-positive} if it is $2$-positive.
\end{definition}

\begin{theorem}[Completely positive implies two-positive]
    \label{thm:cp_implies_two_positive}
    \lean{IsCPMap.is2PositiveMap}
    \leanok
    \uses{def:cp_map, def:k_positive_map}
    Every completely positive map is two-positive.
\end{theorem}

\begin{proof}
    \leanok
    Let $E(X)=\sum_i K_iXK_i^\dagger$ be a Kraus representation. For every
    $Y\geq0$ in $\MN{D}\otimes\MN{2}$,
    \[
        (E\otimes\id_2)(Y)
        =\sum_i (K_i\otimes\Id_2)Y(K_i\otimes\Id_2)^\dagger\geq0,
    \]
    because each summand is positive semidefinite. Thus
    $E\otimes\id_2$ is positive, so $E$ is two-positive.
\end{proof}

\begin{theorem}[Two-positive implies positive]
    \label{thm:two_positive_implies_positive}
    \lean{Is2PositiveMap.isPositiveMap}
    \leanok
    \uses{def:positive_map, def:k_positive_map}
    Every two-positive map is positive.
\end{theorem}

\begin{proof}
    \leanok
    If $X\geq0$, then $X\otimes E_{11}\geq0$ in
    $\MN{D}\otimes\MN{2}$. Two-positivity gives
    \[
        (E\otimes\id_2)(X\otimes E_{11})=E(X)\otimes E_{11}\geq0.
    \]
    Its $(1,1)$ corner is $E(X)$, hence $E(X)\geq0$. Therefore $E$ is
    positive.
\end{proof}

\begin{definition}[Unital linear map]\label{def:unital_linear_map}
    \lean{KadisonSchwarz.IsUnitalMap}
    \leanok
    A linear map $E$ is unital if $E(\Id)=\Id$.
\end{definition}

\begin{theorem}[Kadison--Schwarz for unital two-positive maps]\label{thm:kadison_schwarz_2positive}
    \lean{kadison_schwarz_2positive}
    \leanok
    \uses{def:k_positive_map, def:unital_linear_map}
    If $E$ is unital and two-positive, then for all $X\in\MN{D}$,
    $E(X^\dagger X)\ge E(X)^\dagger E(X)$.
\end{theorem}

\begin{proof}
    \leanok
    Form the $2\times 2$ block matrix
    $(\begin{smallmatrix} \Id & X \\ X^\dagger & X^\dagger X \end{smallmatrix})\ge0$.
    Apply $E\otimes\id_{2}$ (positive by $2$-positivity) and take the Schur
    complement of the top-left identity block:
    $E(X^\dagger X)-E(X)^\dagger E(X)\ge0$.
\end{proof}

\begin{definition}[Blockwise ampliation]\label{def:blockwise_ampliation}
    \lean{nPositiveAmpliation}
    \leanok
    For a linear map
    $E : \MN{D_{\rm in}}\to\MN{D_{\rm out}}$, its $k$-fold ampliation acts
    between the corresponding block-matrix algebras by
    \begin{align}
        E^{(k)}(X)_{(i,p),(j,q)}
         & =E((X_{(a,p),(b,q)})_{a,b})_{i,j},
        \qquad i,j<D_{\rm out},\quad a,b<D_{\rm in}.
        \label{eq:schwarz_blockwise_ampliation}
    \end{align}
\end{definition}

\begin{theorem}[$k$-positivity as positivity of the ampliation]
    \lean{isNPositiveMap_iff_isPositiveMap_nPositiveAmpliation}
    \leanok
    \uses{def:k_positive_map, def:blockwise_ampliation}
    A linear map is $k$-positive if and only if its blockwise ampliation is
    positive on $\MN{D}\otimes\MN{k}$.
\end{theorem}

\begin{proof}
    \leanok
    This is precisely the preceding definition of the ampliation.
\end{proof}

\begin{definition}[Schmidt rank]\label{def:schmidt_rank}
    \lean{Matrix.schmidtCoeffMatrix, Matrix.schmidtRank, Matrix.HasSchmidtRankLE}
    \leanok
    A vector $\psi\in\C^{D}\otimes\C^{k}$ is identified with its coefficient
    matrix $C_\psi\in \MN{D,k}(\C)$. Its Schmidt rank is
    $\SR(\psi)=\rank(C_\psi)$. We write $\SR(\psi)\le r$ for the corresponding
    bounded-rank condition.
\end{definition}

\begin{theorem}[Basic bounds for Schmidt rank]\label{thm:schmidt_rank_basic_bounds}
    \lean{Matrix.schmidtRank_zero, Matrix.schmidtRank_le_left,
        Matrix.schmidtRank_le_right, Matrix.schmidtRank_le_min,
        Matrix.schmidtRank_product_le_one, Matrix.hasSchmidtRankLE_one_product}
    \leanok
    \uses{def:schmidt_rank}
    The zero vector has Schmidt rank zero, every vector satisfies
    $\SR(\psi)\le \min(D,k)$, and every product vector $u\otimes v$ has
    Schmidt rank at most one.
\end{theorem}

\begin{proof}
    \leanok
    The zero-vector claim is $\rank(0)=0$. For the dimension bounds,
    $\SR(\psi)=\rank(C_\psi)\le D$ and
    $\SR(\psi)=\rank(C_\psi)\le k$ follow from the row-rank and
    column-rank bounds on a $D\times k$ matrix, hence
    $\SR(\psi)\le \min(D,k)$. For a product vector, each column
    of $C_{u\otimes v}$ is a scalar multiple of $u$, so the column span has
    dimension at most one.
\end{proof}

\begin{definition}[Schmidt singular values]\label{def:schmidt_singular_values}
    \lean{Matrix.schmidtSingularValues}
    \leanok
    \uses{def:schmidt_rank}
    The Schmidt singular values of a vector
    $\psi\in\C^{D}\otimes\C^k$ are the singular values of the coefficient
    matrix $C_\psi$, viewed as a linear map $\C^k\to\C^D$.
\end{definition}

\begin{theorem}[Schmidt singular values and Schmidt rank]
    \label{thm:schmidt_singular_values_support}
    \lean{Matrix.support_schmidtSingularValues,
        Matrix.schmidtSingularValues_eq_zero_iff,
        Matrix.schmidtSingularValues_pos_iff_lt_schmidtRank,
        Matrix.hasSchmidtRankLE_iff_schmidtSingularValues_eq_zero}
    \leanok
    \uses{def:schmidt_rank, def:schmidt_singular_values}
    The Schmidt singular values are indexed from zero. The nonzero Schmidt
    singular values of $\psi$ are exactly those with index below
    $\SR(\psi)$. Equivalently, for zero-based indexing,
    $s_k(\psi)=0$ if and only if $\SR(\psi)\le k$.
\end{theorem}

\begin{proof}
    \leanok
    For a finite-dimensional linear map, the support of the singular-value
    sequence is the initial interval whose length is the dimension of the
    range. The coefficient matrix $C_\psi$ represents the corresponding map
    $\C^k\to\C^D$, and this range dimension is $\rank(C_\psi)$.
\end{proof}

\subsection{Schmidt rank and maximal overlap}

The monotonicity, rank-factorization, reduced-density, Rayleigh, and Parseval
steps used below are proved in
Section~\ref{sec:app_positive_not_cp_schmidt_rank}. The two extremal
projection bounds behind Ky Fan's maximum principle are proved in
Section~\ref{sec:app_positive_not_cp_ky_fan}.


\begin{definition}[Sum of the largest eigenvalues]
    \label{def:ky_fan_norm}
    \lean{Matrix.IsHermitian.kyFanNorm,
        Matrix.IsHermitian.kyFanNorm_eq_trace_re_of_card_le}
    \leanok
    Let $A\in\MN{D}$ be Hermitian with eigenvalues
    $\lambda_1\ge\lambda_2\ge\cdots\ge\lambda_D$ listed in decreasing order.
    For $0\le k\le D$, set
    \begin{align}
        S_k(A)
         & =\sum_{i=1}^{k}\lambda_i.
        \label{eq:channel_largest_eigenvalue_sum}
    \end{align}
    This is the sum of the $k$ largest signed eigenvalues. Indices beyond $D$
    contribute zero, so the value stabilizes at $\re\tr(A)$ once $k$ reaches
    $D$.

    The \emph{Ky-Fan $k$-norm} $\|A\|_{(k)}$ is by definition the sum of the
    $k$ largest \emph{singular values} of $A$. It agrees with $S_k(A)$ exactly
    when $A$ is positive semidefinite, where eigenvalues and singular values
    coincide; for an indefinite Hermitian $A$ the two differ. For
    $A=\diag(1,-2)$, one has $S_1(A)=1$ but $\|A\|_{(1)}=2$. The results below
    concern $S_k(A)$ for arbitrary Hermitian $A$, so on the
    positive-semidefinite cone they are statements about the Ky-Fan norm.
\end{definition}

\begin{theorem}[Sum of the remaining eigenvalues]
    \label{thm:ky_fan_eigenvalue_tail}
    \lean{Matrix.IsHermitian.sum_eigenvalues₀_filter_le}
    \leanok
    \uses{def:ky_fan_norm}
    Let $A\in\MN{D}$ be Hermitian, with eigenvalues
    $\lambda_1\ge\cdots\ge\lambda_D$. For every nonnegative integer $k$,
    \begin{align}
        \sum_{k<i\le D}\lambda_i
         & =\re\tr(A)-S_k(A).
        \label{eq:ky_fan_eigenvalue_tail}
    \end{align}
    If $k\ge D$, the sum on the left is empty and both sides vanish.
    No positivity assumption on $A$ is required.
\end{theorem}

\begin{proof}\leanok
    \uses{def:ky_fan_norm}
    The sum of all eigenvalues is $\re\tr(A)$. Subtract the first $k$
    eigenvalues, with the sum understood to stop at $D$.
\end{proof}

\begin{theorem}[Ky Fan's maximum principle]
    \label{thm:ky_fan_maximum_principle}
    \lean{Matrix.IsHermitian.kyFanNorm_eq_sup_trace}
    \leanok
    \uses{def:ky_fan_norm, thm:ky_fan_achievability, thm:ky_fan_upper_bound}
    Let $A\in\MN{D}$ be Hermitian and $0\le k<D$. Then the sum of its $k$
    largest eigenvalues is the largest value of $\re\tr(PA)$ attained by an
    orthogonal projection $P$ of rank $k$:
    \begin{align}
        S_k(A)
         & =\max_{\substack{P^2=P=P^\dagger\\ \tr(P)=k}}\re\tr(PA).
        \label{eq:channel_ky_fan_maximum}
    \end{align}
    This is Ky Fan's maximum principle \cite{Fan1949Theorem}; see also
    \cite[Problem~I.6.15, Exercise~II.1.13]{Bhatia1997Matrix}. For a positive
    semidefinite $A$, where $S_k(A)$ coincides with the Ky-Fan norm
    $\|A\|_{(k)}$, it gives the extremal overlap of
    \cite[Lemma~3.1]{Wolf2012Quantum}.
\end{theorem}

\begin{proof}\leanok
    \uses{thm:ky_fan_achievability, thm:ky_fan_upper_bound}
    Theorem~\ref{thm:ky_fan_achievability} provides a rank-$k$ projection
    whose trace against $A$ equals $S_k(A)$, so the value is attained.
    Theorem~\ref{thm:ky_fan_upper_bound} shows that no rank-$k$ projection
    exceeds it. Hence $S_k(A)$ is the maximum.
\end{proof}

\begin{theorem}[Maximal overlap with a fixed Schmidt rank]
    \label{thm:maximal_overlap_schmidt_rank}
    \lean{Matrix.maximalSchmidtOverlap_eq_kyFanNorm}
    \leanok
    \uses{def:schmidt_rank, def:ky_fan_norm, thm:ky_fan_maximum_principle}
    Let $\phi\in\C^{D}\otimes\C^k$ be a normalized vector with reduced density
    matrix $\rho=\tr_2\ketbra{\phi}{\phi}$ on the first factor, and let
    $1\le n<D$. The largest squared overlap $|\braket{\phi}{\psi}|^2$ attained
    by a normalized vector $\psi$ of Schmidt rank at most $n$ is the Ky-Fan
    $n$-norm of $\rho$:
    \begin{align}
        \max_{\substack{\SR(\psi)\le n\\ \|\psi\|=1}}
        |\braket{\phi}{\psi}|^2
         & =\|\rho\|_{(n)}.
        \label{eq:schwarz_maximal_schmidt_overlap}
    \end{align}
    Because $\rho$ is positive semidefinite, this norm is the sum of its $n$
    largest eigenvalues, which coincide there with its singular values. This is
    \cite[Lemma~3.1]{Wolf2012Quantum} for $1\le n<D$; the top index $n=D$,
    where the value is $\tr\rho=1$, is
    Theorem~\ref{thm:maximal_overlap_schmidt_rank_top}.
\end{theorem}

\begin{proof}
    \leanok
    \uses{thm:ky_fan_maximum_principle}
    Write $C$ for the coefficient matrix of $\phi$, so that
    $\rho=CC^\dagger$ and the overlap with a vector $\psi$ of coefficient
    matrix $B$ is the Frobenius pairing
    $\braket{\phi}{\psi}=\tr(C^\dagger B)$; normalization reads
    $\tr(B^\dagger B)=1$ and the Schmidt-rank constraint reads
    $\rank(B)\le n$.

    For the upper bound, let $P$ be the orthogonal projection onto the column
    space of $B$, so $\rank(P)\le n$ and $PB=B$. The Cauchy--Schwarz inequality
    for the Frobenius pairing gives
    \begin{align}
        |\tr(C^\dagger B)|^2
         & =|\tr((PC)^\dagger B)|^2
        \le \tr\!((PC)^\dagger(PC))\tr(B^\dagger B)
        =\tr(P\rho).
        \notag
    \end{align}
    Since $P$ has rank at most $n$ and $\rho$ is positive semidefinite, the
    positive-semidefinite form of Ky Fan's maximum principle
    (Theorem~\ref{thm:ky_fan_maximum_principle}) gives
    $\tr(P\rho)\le \|\rho\|_{(n)}$. Thus
    $|\braket{\phi}{\psi}|^2\le\|\rho\|_{(n)}$ for every admissible $\psi$.

    For the converse, let $P_n$ be the eigenprojection of $\rho$ onto its $n$
    largest eigenvalues, so $\tr(P_n\rho)=\|\rho\|_{(n)}$, and set
    \begin{align}
        B
         & =\|\rho\|_{(n)}^{-1/2}P_nC.
        \label{eq:schwarz_schmidt_overlap_witness}
    \end{align}
    Then $\rank(B)\le\rank(P_n)=n$ and
    $\tr(B^\dagger B)=\|\rho\|_{(n)}^{-1}\tr(P_n\rho)=1$, while
    \begin{align}
        |\tr(C^\dagger B)|^2
         & =\|\rho\|_{(n)}^{-1}|\tr(P_n\rho)|^2
        =\|\rho\|_{(n)}.
        \notag
    \end{align}
    Thus the bound is attained. For a normalized $\phi$ with $1\le n$, one
    has $\|\rho\|_{(n)}\ge\|\rho\|_{(1)}>0$, since $\tr\rho=1$, which
    justifies the normalization in
    (\ref{eq:schwarz_schmidt_overlap_witness}).
\end{proof}

\begin{theorem}[Maximal overlap at the top Schmidt rank]
    \label{thm:maximal_overlap_schmidt_rank_top}
    \lean{Matrix.maximalSchmidtOverlap_eq_kyFanNorm_top}
    \leanok
    \uses{def:schmidt_rank, def:ky_fan_norm}
    At the top index $n=D$, the Schmidt-rank constraint $\SR(\psi)\le D$
    is vacuous, so the largest squared overlap is the Ky-Fan $D$-norm of
    $\rho$:
    \begin{align}
        \max_{\|\psi\|=1}|\braket{\phi}{\psi}|^2
         & =\|\rho\|_{(D)}.
        \label{eq:schwarz_maximal_overlap_top}
    \end{align}
    The maximum is attained at $\psi=\phi$. The Ky-Fan $D$-norm sums all
    eigenvalues, so for a normalized $\phi$ it reduces to
    $\|\rho\|_{(D)}=\tr\rho=1$. Together with
    Theorem~\ref{thm:maximal_overlap_schmidt_rank}, this covers the full range
    $1\le n\le D$ of \cite[Lemma~3.1]{Wolf2012Quantum}.
\end{theorem}

\begin{proof}
    \leanok
    The Ky-Fan $D$-norm sums all eigenvalues of $\rho$, so
    $\|\rho\|_{(D)}=\tr\rho=\|\phi\|^2=1$. The bound $1$ is attained at
    $\psi=\phi$, whose Schmidt rank is at most $D$; for any normalized
    $\psi$, the Cauchy--Schwarz inequality gives
    $|\braket{\phi}{\psi}|^2\le\|\phi\|^2\|\psi\|^2=1$.
\end{proof}

\begin{theorem}[Homogeneous Schmidt overlap bound]
    \label{thm:homogeneous_schmidt_overlap_bound}
    \lean{Matrix.normSq_inner_le_kyFanNorm}
    \leanok
    \uses{def:schmidt_rank, def:ky_fan_norm}
    Let $\Psi,\Phi\in\mathbb C^m\otimes\mathbb C^n$, with
    $\|\Phi\|=1$ and $\SR(\Phi)\le D$, where $D\ge0$ is an integer.
    Set $\rho=\tr_2|\Psi\rangle\langle\Psi|$. Then
    \begin{align}
        |\langle\Psi,\Phi\rangle|^2 & \le\|\rho\|_{(D)}.
        \label{eq:homogeneous_schmidt_overlap_bound}
    \end{align}
    No normalization of $\Psi$ is required. If $D\ge m$, the right-hand
    side is $\tr\rho=\|\Psi\|^2$. This is the homogeneous upper bound
    underlying \cite[Lemma~3.1]{Wolf2012Quantum}.
\end{theorem}

\begin{proof}
    \leanok
    \uses{thm:ky_fan_upper_bound}
    Write $C$ and $B$ for the coefficient matrices of $\Psi$ and $\Phi$.
    The support projection $P$ of $BB^\dagger$ has rank at most $D$.
    Since $PB=B$ and $\tr(B^\dagger B)=1$, the Hilbert--Schmidt
    Cauchy--Schwarz inequality gives
    \begin{align}
        |\tr(C^\dagger B)|^2 & \le\tr(PCC^\dagger).
        \label{eq:homogeneous_schmidt_overlap_support}
    \end{align}
    For $D<m$, applying the Ky Fan variational bound to
    (\ref{eq:homogeneous_schmidt_overlap_support}) proves
    (\ref{eq:homogeneous_schmidt_overlap_bound}).
    For $D\ge m$, ordinary Cauchy--Schwarz gives the bound
    $\|\Psi\|^2=\tr\rho$.
\end{proof}

\begin{definition}[Trimming projection]
    \label{def:trimming_projection}
    \lean{Matrix.IsTrimmingProjection, Matrix.leftApply}
    \leanok
    \uses{def:ky_fan_norm}
    Let $\psi\in\mathbb C^m\otimes\mathbb C^n$ have coefficient matrix $C$
    and reduced density $\rho=CC^\dagger$ on the first factor, and let
    $D\ge0$ be an integer. A trimming projection of $\psi$ at $D$ is an
    orthogonal projection $P$ on $\mathbb C^m$ of rank at most $D$ with
    $\tr(P\rho)=\|\rho\|_{(D)}$. The trimmed vector is
    $\operatorname{trim}_D\psi=(P\otimes I)\psi$, with coefficient matrix
    $PC$; it keeps the Schmidt terms of the $D$ largest Schmidt
    coefficients, following \cite{Landau2013PolynomialTime}. When a Schmidt
    coefficient is degenerate, every choice of kept Schmidt vectors gives a
    trimming projection, and a trimming projection exists for every $D$.
\end{definition}

\begin{theorem}[Eckart--Young theorem]
    \label{thm:eckart_young_trimming}
    \lean{Matrix.IsTrimmingProjection.norm_inner_le_re_inner_normalized,
        Matrix.IsTrimmingProjection.inner_normalized_eq_sqrt}
    \leanok
    \uses{def:trimming_projection, thm:homogeneous_schmidt_overlap_bound}
    Let $P$ be a trimming projection of $\psi$ at $D$, and put
    $\psi'=\operatorname{trim}_D\psi/\|\operatorname{trim}_D\psi\|$, read as
    $0$ when the trimmed vector vanishes. Then
    \begin{align}
        \langle\psi',\psi\rangle & =\|\rho\|_{(D)}^{1/2}\ge
        |\langle w,\psi\rangle|
        \label{eq:eckart_young_trimming}
    \end{align}
    for every unit vector $w$ of Schmidt rank at most $D$. No normalization
    of $\psi$ is required. This is \cite[Lemma~lem:eckart]{Landau2013PolynomialTime}.
\end{theorem}

\begin{proof}
    \leanok
    The overlap $\langle\operatorname{trim}_D\psi,\psi\rangle$ and the
    squared norm $\|\operatorname{trim}_D\psi\|^2$ both equal
    $\tr(C^\dagger PC)=\tr(P\rho)=\|\rho\|_{(D)}$, which gives the equality
    in (\ref{eq:eckart_young_trimming}). The inequality is
    Theorem~\ref{thm:homogeneous_schmidt_overlap_bound}.
\end{proof}

\begin{theorem}[Trimming error against a vector of bounded Schmidt rank]
    \label{thm:trimming_error_bounded_schmidt_rank}
    \lean{Matrix.IsTrimmingProjection.norm_inner_sub_leftApply_le,
        Matrix.IsTrimmingProjection.norm_inner_sub_le_norm_inner_leftApply}
    \leanok
    \uses{def:trimming_projection, def:schmidt_rank, thm:ky_fan_upper_bound}
    Let $P$ be a trimming projection of $u$ at $k$, and let $v$ have
    Schmidt rank at most $D$. Then
    \begin{align}
        \bigl|\langle u-\operatorname{trim}_k u,v\rangle\bigr|
         & \le\|u\|\,\|v\|\sqrt{\frac{D}{k+1}},
        \label{eq:trimming_error_bounded_schmidt_rank}
    \end{align}
    and consequently $|\langle\operatorname{trim}_k u,v\rangle|\ge
    |\langle u,v\rangle|-\|u\|\|v\|\sqrt{D/(k+1)}$. For unit vectors,
    trimming to $k+1\ge D/\varepsilon^2$ terms costs at most $\varepsilon$
    in overlap. This corrects
    \cite[Lemma~l:trim1]{Landau2013PolynomialTime}, whose printed radius
    $D/\varepsilon$ does not suffice; the printed statement is not
    formalized \cite{gap:lvv13_trimming_radius}.
\end{theorem}

\begin{proof}
    \leanok
    Write $C$ and $B$ for the coefficient matrices of $u$ and $v$, put
    $\rho=CC^\dagger$, and let $T$ be the orthogonal projection onto the
    range of $(I-P)B$. Then $T$ has rank at most $D$, $PT=0$, and the
    Cauchy--Schwarz inequality gives
    $|\langle u-\operatorname{trim}_k u,v\rangle|^2\le
        \tr(T\rho)\,\|v\|^2$. Applying Theorem~\ref{thm:ky_fan_upper_bound}
    to the projection $P+T$ of rank at most $k+D$, with eigenvalues
    $\mu_1\ge\mu_2\ge\cdots$ of $\rho$, gives
    \begin{align}
        \tr(T\rho) & \le\sum_{i=k+1}^{k+D}\mu_i\le D\mu_{k+1}
        \le\frac{D}{k+1}\tr\rho=\frac{D}{k+1}\|u\|^2.
        \label{eq:trimming_error_ky_fan}
    \end{align}
    When $k+D$ reaches the dimension, the same bound follows from
    $\tr((P+T)\rho)\le\tr\rho$.
\end{proof}

\begin{theorem}[Trace-maximizing projections commute]
    \label{thm:trace_maximizing_projection_commutes}
    \lean{Matrix.IsHermitian.commute_of_forall_trace_mul_re_le}
    \leanok
    Let $\rho$ be Hermitian and $P$ an orthogonal projection such that
    $\tr(R\rho)\le\tr(P\rho)$ for every orthogonal projection $R$ with
    $\tr R=\tr P$. Then $P\rho=\rho P$.
\end{theorem}

\begin{proof}
    \leanok
    It suffices that $P\rho(I-P)=0$: then $P\rho=P\rho P$ is Hermitian, so
    $\rho P=(P\rho)^\dagger=P\rho$. Suppose $P\rho(I-P)v\ne0$, and put
    $z=(I-P)v$ and $y=P\rho z\ne0$, so that $Py=y$, $Pz=0$ and
    $\langle y,\rho z\rangle=\|y\|^2$. For real $t$ let $w=y+tz$, and let
    $R$ be $P$ with the projection onto $y$ replaced by the projection
    onto $w$. Since $w$ is orthogonal to the range of $P$ minus the line
    through $y$, $R$ is an orthogonal projection with $\tr R=\tr P$, and
    \begin{align}
        \tr(R\rho)-\tr(P\rho) & =
        \frac{t\bigl(2\|y\|^4+t(\|y\|^2\langle z,\rho z\rangle
        -\|z\|^2\langle y,\rho y\rangle)\bigr)}
        {\|y\|^2\,(\|y\|^2+t^2\|z\|^2)}.
        \label{eq:trace_maximizing_projection_rotation}
    \end{align}
    This is positive for small $t>0$, a contradiction.
\end{proof}

\begin{theorem}[Trimming as a projection on the second factor]
    \label{thm:trimming_second_factor}
    \lean{Matrix.IsTrimmingProjection.norm_leftApply_le,
        Matrix.IsTrimmingProjection.commute_leftMarginal,
        Matrix.rightApply,
        Matrix.IsTrimmingProjection.exists_mul_eq_mul_transpose,
        Matrix.IsTrimmingProjection.exists_leftApply_eq_rightApply}
    \leanok
    \uses{def:trimming_projection, thm:ky_fan_upper_bound,
        thm:trace_maximizing_projection_commutes}
    Let $P$ be a trimming projection of $\psi$ at $D$, with coefficient
    matrix $C$ and $\rho=CC^\dagger$. Then
    $\|\operatorname{trim}_D\psi\|\le\|\psi\|$, the projection $P$ commutes
    with $\rho$, and there is an orthogonal projection $Q$ on
    $\mathbb C^n$ of rank at most $\operatorname{rank}P\le D$ with
    $PC=CQ^{\mathsf T}$, that is
    \begin{align}
        \operatorname{trim}_D\psi & =(P\otimes I)\psi=(I\otimes Q)\psi.
        \label{eq:trimming_second_factor}
    \end{align}
    This is the form of the trimmed vector used in the proof of
    \cite[Lemma~l:trim2]{Landau2013PolynomialTime}.
\end{theorem}

\begin{proof}
    \leanok
    The squared norm of the trimmed vector is $\|\rho\|_{(D)}\le\tr\rho
        =\|\psi\|^2$. Every orthogonal projection $R$ with $\tr R=\tr P$ has
    rank $\operatorname{rank}P\le D$, so
    $\tr(R\rho)\le\|\rho\|_{(D)}=\tr(P\rho)$ by
    Theorem~\ref{thm:ky_fan_upper_bound}, and $P$ commutes with $\rho$ by
    Theorem~\ref{thm:trace_maximizing_projection_commutes}. Let
    $Q^{\mathsf T}$ be the support projection of $C^\dagger PC$, whose rank
    is $\operatorname{rank}(PC)\le\operatorname{rank}P$. It fixes the rows
    of $PC$, so $PCQ^{\mathsf T}=PC$. Its range lies in the range of
    $C^\dagger P$, and $(I-P)CC^\dagger P=(I-P)P\rho=0$, so
    $(I-P)CQ^{\mathsf T}=0$. Adding the two identities gives
    $CQ^{\mathsf T}=PC$.
\end{proof}

\begin{theorem}[Trimming through an isometry on the second factor]
    \label{thm:trimming_isometric_factorization}
    \lean{Matrix.exists_isometry_mul_conjTranspose_eq,
        Matrix.IsTrimmingProjection.exists_isometry_leftApply_eq_rightApply}
    \leanok
    \uses{thm:trimming_second_factor}
    Every orthogonal projection $Q$ on $\mathbb C^n$ is $Q=WW^\dagger$ for an
    isometry $W\colon\mathbb C^k\to\mathbb C^n$, $W^\dagger W=I$, with
    $k=\operatorname{rank}Q$. Consequently, for a trimming projection of
    $\psi$ at $D$ there is an isometry $W\colon\mathbb C^k\to\mathbb C^n$
    with $k\le\min(D,n)$ and
    \begin{align}
        \operatorname{trim}_D\psi & =(I\otimes W)\,x,
                                  & x & =(I\otimes W^\dagger)\psi.
        \label{eq:trimming_isometric_factorization}
    \end{align}
    The isometry sends the $j$-th basis vector to the $j$-th kept right
    Schmidt vector, and $x$ is the left state of the trimmed vector, as in
    the boundary contractions of \cite{Landau2013PolynomialTime}.
\end{theorem}

\begin{proof}
    \leanok
    Diagonalize $Q$ by a unitary; its eigenvalues are $0$ and $1$, and the
    eigenvectors of eigenvalue $1$ form the columns of $W$, so $W^\dagger
        W=I$ and $WW^\dagger=Q$. The rank of $Q=WW^\dagger$ is the rank of
    $W^\dagger W=I_k$, which is $k$. For the trimmed vector take $Q$ from
    Theorem~\ref{thm:trimming_second_factor}: the coefficient matrix of
    $(I\otimes W)(I\otimes W^\dagger)\psi$ is
    $C\,\overline{W}\,W^{\mathsf T}=C(WW^\dagger)^{\mathsf T}=CQ^{\mathsf T}=PC$.
\end{proof}

\begin{theorem}[No cross terms after trimming]
    \label{thm:trimming_no_cross_terms}
    \lean{Matrix.toEuclideanLin_kronecker_one_apply,
        Matrix.toEuclideanLin_one_kronecker_apply,
        Matrix.IsTrimmingProjection.norm_leftApply_sq_add_norm_sub_sq,
        Matrix.IsTrimmingProjection.inner_toEuclideanLin_eq_add}
    \leanok
    \uses{thm:trimming_second_factor}
    Let $P$ be a trimming projection of $\psi$ at $D$, and write
    $w_0=\operatorname{trim}_D\psi$ and $w_1=\psi-w_0$. Then $w_0$ and $w_1$
    are orthogonal, so $\|w_0\|^2+\|w_1\|^2=\|\psi\|^2$, and for every
    operator $X=A\otimes I+I\otimes B$ that is a sum of an operator on each
    factor,
    \begin{align}
        \langle\psi,X\psi\rangle & =\langle w_0,Xw_0\rangle
        +\langle w_1,Xw_1\rangle.
        \label{eq:trimming_no_cross_terms}
    \end{align}
\end{theorem}

\begin{proof}
    \leanok
    Write $\rho=CC^\dagger$. The cross term of $A\otimes I$ is
    $\tr(PA(I-P)\rho)$, which vanishes because $P$ commutes with $\rho$ and
    $(I-P)P=0$; the other cross term of $A\otimes I$ is $\tr((I-P)AP\rho)=0$
    for the same reason. The cross terms of $I\otimes B$ are
    $\tr(C^\dagger P(I-P)CB^{\mathsf T})$ and
    $\tr(C^\dagger(I-P)PCB^{\mathsf T})$, both zero, and the overlap of
    $w_0$ with $w_1$ is $\tr(P(I-P)\rho)=0$.
\end{proof}

\begin{theorem}[Energy after trimming]
    \label{thm:trimming_energy}
    \lean{Matrix.PosSemidef.re_inner_toEuclideanLin_sq_le,
        Matrix.PosSemidef.abs_re_inner_toEuclideanLin_le,
        Matrix.PosSemidef.mul_sq_norm_le_re_inner_toEuclideanLin,
        Matrix.IsTrimmingProjection.re_inner_leftApply_le,
        Matrix.IsTrimmingProjection.re_inner_normalized_le,
        Matrix.IsTrimmingProjection.re_inner_normalized_le_of_le_half}
    \leanok
    \uses{thm:trimming_no_cross_terms}
    Let $H=A\otimes I+I\otimes B+K$ with $H\ge E_0$ and $0\le K\le I$, let
    $\psi$ be a unit vector with $\operatorname{Re}\langle\psi,H\psi\rangle
        \le E_0+\delta$, let $P$ be a trimming projection of $\psi$ at $D$, and
    write $w_0=\operatorname{trim}_D\psi$, $w_1=\psi-w_0$ and
    $t=\|w_1\|^2$. Then
    \begin{align}
        \operatorname{Re}\langle w_0,Hw_0\rangle & \le E_0(1-t)+\delta+2\sqrt t,
        \label{eq:trimming_energy_unnormalized}
    \end{align}
    and if $t<1$ the normalized trimmed vector $v=w_0/\|w_0\|$ satisfies
    \begin{align}
        \operatorname{Re}\langle v,Hv\rangle & \le
        E_0+\frac{\delta+2\sqrt t}{1-t},
        \label{eq:trimming_energy_normalized}
    \end{align}
    which is at most $E_0+2(\delta+2\sqrt t)$ when $t\le\frac12$. This is the
    corrected core of \cite[Lemma~l:2]{Landau2013PolynomialTime}; see
    \cite{gap:lvv13_trimming_energy}.
\end{theorem}

\begin{proof}
    \leanok
    By Theorem~\ref{thm:trimming_no_cross_terms} applied to
    $X=A\otimes I+I\otimes B$ and the expansion of $\langle\psi,K\psi\rangle$
    along $\psi=w_0+w_1$,
    $\langle w_0,Hw_0\rangle=\langle\psi,H\psi\rangle-\langle w_1,Hw_1\rangle
        -2\operatorname{Re}\langle w_0,Kw_1\rangle$. Here
    $\langle w_1,Hw_1\rangle\ge E_0t$. The Cauchy--Schwarz inequality for the
    positive form of $K$, obtained from the discriminant of
    $s\mapsto\langle w_0+sw_1,K(w_0+sw_1)\rangle\ge0$, together with $K\le I$,
    gives $|\operatorname{Re}\langle w_0,Kw_1\rangle|\le\|w_0\|\|w_1\|\le\sqrt t$.
    This proves (\ref{eq:trimming_energy_unnormalized}). Dividing by
    $\|w_0\|^2=1-t$ gives (\ref{eq:trimming_energy_normalized}); since the
    energy bound forces $\delta\ge0$, the last bound follows from
    $1/(1-t)\le2$.
\end{proof}

\begin{theorem}[Trimmed mass from an overlap]
    \label{thm:trimming_mass}
    \lean{Matrix.IsTrimmingProjection.norm_sub_leftApply_sq_le}
    \leanok
    \uses{thm:eckart_young_trimming, thm:trimming_no_cross_terms}
    Let $\psi$ be a unit vector with trimming projection $P$ at $D$, and let
    $u$ be a unit vector of Schmidt rank at most $D$ with
    $|\langle u,\psi\rangle|\ge1-\eta$. Then
    $\|\psi-\operatorname{trim}_D\psi\|^2\le2\eta$.
\end{theorem}

\begin{proof}
    \leanok
    By Theorem~\ref{thm:eckart_young_trimming},
    $\|\operatorname{trim}_D\psi\|\ge|\langle u,\psi\rangle|\ge1-\eta$. If
    $\eta\le1$, squaring gives $\|\operatorname{trim}_D\psi\|^2\ge
        (1-\eta)^2\ge1-2\eta$, and the orthogonality of
    Theorem~\ref{thm:trimming_no_cross_terms} gives the claim. If $\eta>1$
    the claim holds since the left side is at most $1$.
\end{proof}

\begin{theorem}[Energy after trimming, from an overlap]
    \label{thm:trimming_energy_overlap}
    \lean{Matrix.IsTrimmingProjection.re_inner_normalized_le_of_norm_inner}
    \leanok
    \uses{thm:trimming_energy, thm:trimming_mass}
    Under the hypotheses of Theorem~\ref{thm:trimming_energy}, let $u$ be a
    unit vector of Schmidt rank at most $D$ with
    $|\langle u,\psi\rangle|\ge1-\eta$ and $2\eta<1$. Then the normalized
    trimmed vector $v$ satisfies
    \begin{align}
        \operatorname{Re}\langle v,Hv\rangle & \le
        E_0+\frac{\delta+2\sqrt{2\eta}}{1-2\eta}.
        \label{eq:trimming_energy_overlap}
    \end{align}
    In \cite[Lemma~l:2]{Landau2013PolynomialTime}, $u$ is the normalized
    trimmed ground state and Theorem~\ref{thm:overlap_through_common_vector}
    gives $\eta=2\delta(1+1/\varepsilon)$, where $\varepsilon$ is the
    spectral gap. The resulting bound depends on $1/\varepsilon$; the bound
    $E_0+6\sqrt\delta$ stated there does not follow from its argument and is
    not formalized here \cite{gap:lvv13_trimming_energy}.
\end{theorem}

\begin{proof}
    \leanok
    Theorem~\ref{thm:trimming_mass} gives $t\le2\eta<1$, so
    (\ref{eq:trimming_energy_normalized}) applies, and its right side
    increases with $t$ because $\delta\ge0$.
\end{proof}

\begin{theorem}[A trimmed low-energy witness on a subspace]
    \label{thm:trimmed_witness}
    \lean{Matrix.ofLp_rightApply, Matrix.leftApply_smul, Matrix.leftApply_rightApply,
        Matrix.schmidtRank_leftApply_le_rank,
        Matrix.IsTrimmingProjection.leftApply_leftApply_of_leftApply_eq,
        Matrix.IsTrimmingProjection.exists_trimmed_witness}
    \leanok
    \uses{thm:trimming_energy_overlap, thm:trimming_isometric_factorization}
    Under the hypotheses of Theorem~\ref{thm:trimming_energy_overlap}, let the
    left factor of $\psi$ lie in a subspace $V$, that is
    $(\Pi_V\otimes1)\psi=\psi$. Then there is a unit vector $w$ of Schmidt rank
    at most $D$ with
    $\operatorname{Re}\langle w,Hw\rangle\le E_0+(\delta+2\sqrt{2\eta})/(1-2\eta)$
    and $(\Pi_V\otimes1)w=w$, together with an isometry
    $W:\mathbb C^k\to\mathbb C^n$, $k\le D$, and a vector $x$ with
    $(\Pi_V\otimes1)x=x$ and $w=(1\otimes W)x$. This is
    Lemma~l:within of~\cite[lines 509--515]{Landau2013PolynomialTime}, with
    the corrected energy bound; the energy $\varepsilon_0+6\sqrt{c_\varepsilon}$
    printed there is not established by its argument
    \cite{gap:lvv13_trimming_energy}.
\end{theorem}

\begin{proof}
    \leanok
    Let $w$ be the normalized trimming of $\psi$ to $D$ terms; it is nonzero
    because the discarded part has squared norm at most $2\eta<1$. Its energy is
    Theorem~\ref{thm:trimming_energy_overlap}. The trimmed vector is
    $(1\otimes Q)\psi$ for a projection $Q$, which commutes with
    $\Pi_V\otimes1$, so it stays in $V\otimes\mathbb C^n$.
    Theorem~\ref{thm:trimming_isometric_factorization} writes it as
    $(1\otimes W)(1\otimes W^\dagger)\psi$, and
    $x=\lVert\cdot\rVert^{-1}(1\otimes W^\dagger)\psi$ is again in
    $V\otimes\mathbb C^k$.
\end{proof}

\begin{definition}[Reduced density of an eigenvector]
    \label{def:reduced_eig_density}
    \lean{Matrix.IsHermitian.reducedEigDensity}
    \leanok
    Let $\tau$ be a Hermitian operator on
    $\C^{D'}\otimes\C^{D}$ with normalized eigenvectors $\phi_i$. The reduced
    density operator of the $i$-th eigenvector on the first factor is
    $\rho_i=\tr_2\ketbra{\phi_i}{\phi_i}$.
\end{definition}

\begin{theorem}[Spectral lower bound for the Schmidt-rank expectation]
    \label{thm:spectral_lower_bound_schmidt_rank}
    \lean{Matrix.IsHermitian.spectral_lower_bound}
    \leanok
    \uses{def:schmidt_rank, def:ky_fan_norm, def:reduced_eig_density,
        lem:reduced_eig_density_posSemidef, thm:maximal_overlap_schmidt_rank}
    Let $\tau$ be a Hermitian operator on $\C^{D'}\otimes\C^{D}$ with
    eigenvalues $\nu_i$ and normalized eigenvectors $\phi_i$, and write
    $\rho_i=\tr_2\ketbra{\phi_i}{\phi_i}$ for the reduced density operator of
    $\phi_i$ on the first factor. Let $1\le n<D'$ and let $\nu_0\ge 0$ be a
    lower bound for the positive eigenvalues (the smallest positive
    eigenvalue). Then every normalized vector $\psi$ of Schmidt rank at most
    $n$ satisfies
    \begin{align}
        \nu_0+\sum_{i:\nu_i\le0}
        (\nu_i-\nu_0)\|\rho_i\|_{(n)}
         & \le\bra{\psi}\tau\ket{\psi}.
        \label{eq:schwarz_spectral_schmidt_lower}
    \end{align}
    Hence the same bound holds for the infimum over such vectors. This is the
    lower bound of \cite[Chapter~3, Proposition~3.2]{Wolf2012Quantum} for
    $n<D'$; the top index $n=D'$ is
    Theorem~\ref{thm:spectral_bound_schmidt_rank_top}.
\end{theorem}

\begin{proof}
    \leanok
    \uses{lem:eigenbasis_rayleigh, lem:eigenbasis_parseval,
        thm:maximal_overlap_schmidt_rank}
    By (\ref{eq:schwarz_eigenbasis_rayleigh}),
    $\bra{\psi}\tau\ket{\psi}
        =\sum_i\nu_i|\braket{\phi_i}{\psi}|^2$, while
    (\ref{eq:schwarz_eigenbasis_parseval}) gives
    $\sum_i|\braket{\phi_i}{\psi}|^2=\|\psi\|^2=1$. Subtracting $\nu_0$
    times the latter from the former gives
    \begin{align}
        \bra{\psi}\tau\ket{\psi}
         & =\nu_0+\sum_i(\nu_i-\nu_0)
        |\braket{\phi_i}{\psi}|^2.
        \label{eq:schwarz_spectral_lower_expansion}
    \end{align}
    For an eigenvalue $\nu_i>0$, the factor $\nu_i-\nu_0$ is non-negative, so
    those terms only increase the sum and may be dropped. For an eigenvalue
    $\nu_i\le0$, the factor $\nu_i-\nu_0$ is nonpositive, and
    Theorem~\ref{thm:maximal_overlap_schmidt_rank} gives
    $|\braket{\phi_i}{\psi}|^2\le\|\rho_i\|_{(n)}$. Multiplication by the
    nonpositive factor reverses the inequality, which gives
    (\ref{eq:schwarz_spectral_schmidt_lower}).
\end{proof}

\begin{theorem}[Spectral upper bound for the Schmidt-rank expectation]
    \label{thm:spectral_upper_bound_schmidt_rank}
    \lean{Matrix.IsHermitian.exists_le_spectral_upper_bound}
    \leanok
    \uses{def:schmidt_rank, def:ky_fan_norm, def:reduced_eig_density,
        lem:reduced_eig_density_posSemidef, thm:maximal_overlap_schmidt_rank}
    With $\tau$, $\phi_i$, and $\rho_i$ as above, let $1\le n<D'$ and suppose
    every eigenvalue of $\tau$ is at most $\nu$. For any eigenvector index $j$,
    there is a normalized vector $\psi$ of Schmidt rank at most $n$ with
    \begin{align}
        \bra{\psi}\tau\ket{\psi}
         & \le\nu+(\nu_j-\nu)\|\rho_j\|_{(n)}.
        \label{eq:schwarz_spectral_schmidt_upper}
    \end{align}
    Hence the infimum over such vectors lies below that value. Taking $\nu$
    to be the largest positive eigenvalue and $j$ the index of the unique
    non-positive eigenvalue $\nu_-$ recovers the upper bound
    $\nu+(\nu_--\nu)\|\rho_-\|_{(n)}$ of
    \cite[Chapter~3, Proposition~3.2]{Wolf2012Quantum}. The statement is for
    $n<D'$; the top index $n=D'$ is
    Theorem~\ref{thm:spectral_bound_schmidt_rank_top}.
\end{theorem}

\begin{proof}
    \leanok
    \uses{lem:eigenbasis_rayleigh, lem:eigenbasis_parseval,
        thm:maximal_overlap_schmidt_rank}
    Replacing $\nu_0$ by $\nu$ in
    (\ref{eq:schwarz_spectral_lower_expansion}) gives
    \begin{align}
        \bra{\psi}\tau\ket{\psi}
         & =\nu+\sum_i(\nu_i-\nu)|\braket{\phi_i}{\psi}|^2.
        \notag
    \end{align}
    Every factor $\nu_i-\nu$ is nonpositive, so all terms are nonpositive and
    keeping only the $i=j$ term gives
    $\bra{\psi}\tau\ket{\psi}
        \le\nu+(\nu_j-\nu)|\braket{\phi_j}{\psi}|^2$.
    Theorem~\ref{thm:maximal_overlap_schmidt_rank} supplies a normalized
    $\psi$ of Schmidt rank at most $n$ attaining
    $|\braket{\phi_j}{\psi}|^2=\|\rho_j\|_{(n)}$; at that vector the
    right-hand side is the expression in
    (\ref{eq:schwarz_spectral_schmidt_upper}).
\end{proof}

\begin{theorem}[Spectral criterion at the top Schmidt rank]
    \label{thm:spectral_bound_schmidt_rank_top}
    \lean{Matrix.IsHermitian.spectral_lower_bound_top,
        Matrix.IsHermitian.exists_spectral_lower_bound_top}
    \leanok
    \uses{lem:eigenbasis_rayleigh, lem:eigenbasis_parseval}
    At the top index $n=D'$, the Schmidt-rank constraint is vacuous and each
    Ky-Fan $D'$-norm reduces to $\|\rho_i\|_{(D')}=\tr\rho_i=1$. The Rayleigh
    characterization of the least eigenvalue then applies: every normalized
    $\psi$ satisfies $\nu_{\min}\le\bra{\psi}\tau\ket{\psi}$, and a
    least-eigenvalue eigenvector attains
    $\bra{\psi}\tau\ket{\psi}=\nu_{\min}$. Together they give
    \begin{align}
        \inf_{\|\psi\|=1}\bra{\psi}\tau\ket{\psi}
         & =\nu_{\min}.
        \label{eq:schwarz_spectral_schmidt_top}
    \end{align}
    There is no Schmidt-rank restriction. At $n=D'$ the right-hand side of the
    upper bound of \cite[Chapter~3, Proposition~3.2]{Wolf2012Quantum} reads
    $\nu+(\nu_--\nu)\|\rho_-\|_{(D')}=\nu_-=\nu_{\min}$, so the display is that
    bound at that index. The positivity index of
    \cite[Chapter~3, Proposition~3.1]{Wolf2012Quantum} runs over $1\le n\le D$
    with $D$ the dimension of the second factor, so $n=D'$ is an index the
    source speaks about when $D'\le D$, and lies beyond its range when $D'>D$.
    The lower bound of the same proposition reads, at that index,
    $\nu_0+\sum_{i:\nu_i\le0}(\nu_i-\nu_0)$, an inequality no stronger than
    (\ref{eq:schwarz_spectral_schmidt_top}), and
    Theorem~\ref{thm:spectral_infimum_top} establishes it. With
    Theorems~\ref{thm:spectral_lower_bound_schmidt_rank}
    and~\ref{thm:spectral_upper_bound_schmidt_rank}, both bounds then hold at
    every $n\ge1$.
\end{theorem}

\begin{proof}
    \leanok
    \uses{lem:eigenbasis_rayleigh, lem:eigenbasis_parseval}
    The eigenbasis expansion (\ref{eq:schwarz_eigenbasis_rayleigh}), together
    with Parseval's identity (\ref{eq:schwarz_eigenbasis_parseval}), bounds
    each weight $\nu_i$ below by $\nu_{\min}$, giving
    $\nu_{\min}\le\bra{\psi}\tau\ket{\psi}$. Evaluating at the eigenvector
    $\phi_j$ for a least-eigenvalue index $j$ gives
    $\bra{\phi_j}\tau\ket{\phi_j}=\nu_j=\nu_{\min}$, so the lower bound is
    attained and the infimum is (\ref{eq:schwarz_spectral_schmidt_top}).
\end{proof}

\begin{theorem}[Rank-one test for $k$-positivity]\label{thm:k_positive_rank_one_test}
    \lean{isNPositiveMap_iff_forall_ampliation_rank_one_posSemidef}
    \leanok
    \uses{def:k_positive_map, def:blockwise_ampliation}
    A linear map $E:\MN{d}\to\MN{d'}$ is $k$-positive if and only if, for
    every vector $\phi\in\C^d\otimes\C^k$,
    $E^{(k)}(\ketbra{\phi}{\phi})\ge0$.
    This is the pure-state reduction used in
    \cite[Chapter~3, Proposition~3.1]{Wolf2012Quantum}.
\end{theorem}

\begin{proof}
    \leanok
    The forward implication is the definition of $k$-positivity, since
    $\ketbra{\phi}{\phi}$ is positive semidefinite. Conversely, every
    positive semidefinite matrix on $\C^d\otimes\C^k$ is a finite sum
    $\sum_i\ketbra{\phi_i}{\phi_i}$, and
    \begin{align}
        E^{(k)}\!\left(\sum_i\ketbra{\phi_i}{\phi_i}\right)
         & =\sum_iE^{(k)}(\ketbra{\phi_i}{\phi_i})\ge0.
        \notag
    \end{align}
    by linearity of $E^{(k)}$, the hypothesis, and closure of the positive
    semidefinite cone under finite sums.
\end{proof}

\subsection{Choi compression criteria}

The right-tensor identities, parametrizations of bounded Schmidt rank, and
projection factorizations used in this subsection are proved in
Section~\ref{sec:app_positive_not_cp_choi_compression}.
Wolf states item~3 of Proposition~3.1 using Schmidt rank exactly $k$.
Under the stated assumption $k\le d$, that test set is empty when $k>d'$.
The dimension-generic criterion below therefore uses Schmidt rank at most
$k$; the source wording is analyzed in
\cite{gap:wolf_prop3_1_exact_schmidt_rank_scope}.

\begin{definition}[Right compression of the Choi matrix]\label{def:choi_right_compression}
    \lean{ChoiJamiolkowski.rightCompression,
        ChoiJamiolkowski.rightCompression_apply,
        ChoiJamiolkowski.compressedOmegaVector}
    \leanok
    \uses{def:blockwise_ampliation}
    Let $T:\MN{d}\to\MN{d'}$ have Choi matrix $\tau$ on
    $\C^{d'}\otimes\C^d$. For $X\in M_{d\times k}(\C)$, the right-factor
    compression is the matrix on $\C^{d'}\otimes\C^k$ with entries
    \begin{align}
        C_{(i,p),(j,q)}
         & =\sum_{a,b}X_{a,p}\tau_{(i,a),(j,b)}
        \overline{X_{b,q}}.
        \label{eq:schwarz_right_choi_compression}
    \end{align}
    The associated vector on $\C^d\otimes\C^k$ has coefficients
    $d^{-1/2}X_{a,p}$.
\end{definition}

\begin{definition}[Right tensor factor for Choi compression]
    \label{def:choi_right_tensor_factor}
    \lean{ChoiJamiolkowski.rightTensorMatrix}
    \leanok
    For $X\in M_{d\times k}(\C)$, let $R_X$ be the matrix from
    $\C^{d'}\otimes\C^d$ to $\C^{d'}\otimes\C^k$ with entries
    \begin{align}
        (R_X)_{(i,p),(j,a)}
         & =\delta_{i,j}X_{a,p}.
        \label{eq:schwarz_right_tensor_factor}
    \end{align}
\end{definition}
\begin{theorem}[Rectangular Schmidt-rank Choi expectation from $k$-positivity]
    \label{thm:k_positive_schmidt_rank_choi_expectation}
    \lean{ChoiJamiolkowski.IsNPositiveMap.choiMatrix_quadraticForm_nonneg_of_hasSchmidtRankLE_rectangular}
    \leanok
    \uses{lem:right_tensor_adjoint_schmidt_rank_representatives,
        thm:choi_right_compression_quadratic_form,
        thm:k_positive_iff_right_choi_compressions}
    Assume $d>0$. If $T:\MN{d}\to\MN{d'}$ is $k$-positive, then
    $\langle\psi,\tau_T\psi\rangle\ge0$ for every
    $\psi\in\C^{d'}\otimes\C^d$ with $\SR(\psi)\le k$.
\end{theorem}

\begin{proof}
    \leanok
    Write $\psi=R_X^\dagger\eta$ using
    Lemma~\ref{lem:right_tensor_adjoint_schmidt_rank_representatives}. By
    Theorem~\ref{thm:k_positive_iff_right_choi_compressions}, the
    compression $C_T(X)$ is positive semidefinite. Its quadratic form at
    $\eta$ is therefore non-negative, and
    (\ref{eq:schwarz_right_compression_quadratic}) identifies this number
    with $\langle\psi,\tau_T\psi\rangle$.
\end{proof}

\begin{theorem}[Converse rectangular Schmidt-rank Choi test]
    \label{thm:schmidt_rank_choi_test_converse}
    \lean{ChoiJamiolkowski.isNPositiveMap_of_forall_hasSchmidtRankLE_choiMatrix_quadraticForm_nonneg_rectangular}
    \leanok
    \uses{lem:right_tensor_adjoint_schmidt_rank,
        thm:choi_right_compression_quadratic_form,
        thm:k_positive_iff_right_choi_compressions}
    Assume $d>0$. Let $T:\MN{d}\to\MN{d'}$. If
    $\langle\psi,\tau_T\psi\rangle\ge0$ for every
    $\psi\in\C^{d'}\otimes\C^d$ with $\SR(\psi)\le k$, then $T$ is
    $k$-positive. This includes $k=0$ and imposes no nonemptiness assumption
    on the output factor.
\end{theorem}

\begin{proof}
    \leanok
    By Theorem~\ref{thm:k_positive_iff_right_choi_compressions}, it is enough
    to prove that $C_T(X)$ is positive semidefinite for each
    $X\in M_{d\times k}(\C)$. For any $\eta$,
    (\ref{eq:schwarz_right_compression_quadratic}) identifies
    $\langle\eta,C_T(X)\eta\rangle$ with
    $\langle R_X^\dagger\eta,\tau_T R_X^\dagger\eta\rangle$, and
    Lemma~\ref{lem:right_tensor_adjoint_schmidt_rank} gives
    $\SR(R_X^\dagger\eta)\le k$. Thus every quadratic form of $C_T(X)$ is
    non-negative. Over a complex finite-dimensional space, polarization
    recovers Hermiticity from this real non-negative quadratic-form condition,
    so $C_T(X)$ is positive semidefinite.
\end{proof}

\begin{theorem}[Rectangular Schmidt-rank Choi criterion]
    \label{thm:schmidt_rank_choi_test_iff}
    \lean{ChoiJamiolkowski.isNPositiveMap_iff_forall_hasSchmidtRankLE_choiMatrix_quadraticForm_nonneg_rectangular}
    \leanok
    \uses{thm:k_positive_schmidt_rank_choi_expectation,
        thm:schmidt_rank_choi_test_converse}
    Assume $d>0$. A map $T:\MN{d}\to\MN{d'}$ is $k$-positive if and only if
    $\langle\psi,\tau_T\psi\rangle\ge0$ for every
    $\psi\in\C^{d'}\otimes\C^d$ with $\SR(\psi)\le k$.
\end{theorem}

\begin{proof}
    \leanok
    Combine Theorem~\ref{thm:k_positive_schmidt_rank_choi_expectation} with
    Theorem~\ref{thm:schmidt_rank_choi_test_converse}.
\end{proof}
\begin{theorem}[Rectangular Choi-compression test for $k$-positivity]
    \label{thm:k_positive_iff_right_choi_compressions}
    \lean{ChoiJamiolkowski.isNPositiveMap_iff_forall_rightCompression_posSemidef}
    \leanok
    \uses{thm:k_positive_rank_one_test, def:choi_right_compression,
        lem:rank_one_ampliation_choi_compression}
    Assume $d>0$. A linear map
    $T:\MN{d}\to\MN{d'}$ is $k$-positive if and only if, for every
    $X\in M_{d\times k}(\C)$, the right-factor compression of the Choi matrix of
    $T$ by $X$ is positive semidefinite.
\end{theorem}

\begin{proof}
    \leanok
    The forward direction applies the rank-one $k$-positivity test to the
    vector $\psi_X$ in (\ref{eq:schwarz_compressed_omega}), then uses
    Lemma~\ref{lem:rank_one_ampliation_choi_compression} in the form
    \begin{align}
        (T\otimes\id_k)(\ketbra{\psi_X}{\psi_X})
         & =C_T(X).
        \label{eq:schwarz_ampliation_compression}
    \end{align}
    Conversely, given a vector $\psi$ on the ampliated space, take
    $X_{a,p}=d^{1/2}\psi_{(a,p)}$. Since $d>0$, this choice satisfies
    $d^{-1/2}X_{a,p}=\psi_{(a,p)}$. Substituting this vector in
    (\ref{eq:schwarz_ampliation_compression}) gives
    $(T\otimes\id_k)(\ketbra{\psi}{\psi})=C_T(X)$.
    Positivity of $C_T(X)$ is therefore positivity of the ampliation on the
    rank-one matrix $\ketbra{\psi}{\psi}$, and the rank-one test gives
    $k$-positivity.
\end{proof}

\begin{theorem}[Right tensor Choi sandwiches of a $k$-positive map]
    \label{thm:k_positive_right_tensor_choi_sandwiches}
    \lean{ChoiJamiolkowski.IsNPositiveMap.rightTensor_choiMatrix_sandwich_posSemidef_rectangular}
    \leanok
    \uses{thm:k_positive_iff_right_choi_compressions,
        thm:choi_right_tensor_sandwich}
    Assume $d>0$. If
    $T:\MN{d}\to\MN{d'}$ is $k$-positive, then, for every
    $X\in M_{d\times k}(\C)$,
    \begin{align}
        R_X\tau R_X^\dagger
         & \ge0,
        \label{eq:schwarz_k_positive_choi_sandwich}
    \end{align}
    where $\tau$ is the Choi matrix of $T$.
\end{theorem}

\begin{proof}
    \leanok
    By Theorem~\ref{thm:k_positive_iff_right_choi_compressions},
    $C_T(X)\ge0$ for every $X\in M_{d\times k}(\C)$. The sandwich formula
    (\ref{eq:schwarz_right_tensor_sandwich}) identifies
    $R_X\tau R_X^\dagger$ with $C_T(X)$.
\end{proof}
\begin{theorem}[Rank-$k$ input projections in the rectangular Choi criterion]
    \label{thm:rank_k_right_projection_choi_sandwich_forward}
    \lean{ChoiJamiolkowski.IsNPositiveMap.rightProjection_choiMatrix_sandwich_posSemidef_rectangular}
    \leanok
    \uses{lem:rank_k_hermitian_projection_factorization,
        thm:factorized_right_projection_choi_sandwich}
    Assume $d>0$. Let $T:\MN{d}\to\MN{d'}$ be $k$-positive,
    with Choi matrix $\tau$. If $P\in\MN{d}$ is Hermitian, satisfies
    $P^2=P$, and has rank $k$, then $R_P\tau R_P^\dagger\geq0$.
    This is the forward rank-$k$ projection-compression direction of
    \cite[Chapter~3, Proposition~3.1]{Wolf2012Quantum}.
\end{theorem}

\begin{proof}
    \leanok
    By Lemma~\ref{lem:rank_k_hermitian_projection_factorization}, write
    $P=VV^\dagger$ with $V\in M_{d\times k}(\C)$. The conclusion is then exactly
    Theorem~\ref{thm:factorized_right_projection_choi_sandwich}.
\end{proof}

\begin{theorem}[Projection-compression converse for the rectangular Choi criterion]
    \label{thm:rank_k_right_projection_choi_sandwich_converse}
    \lean{ChoiJamiolkowski.isNPositiveMap_of_forall_rankProjection_choiMatrix_sandwich_posSemidef_rectangular}
    \leanok
    \uses{thm:k_positive_iff_right_choi_compressions,
        lem:square_sandwich_fixed_rectangular_compression,
        lem:rank_k_projection_fixing_rectangular_factor}
    Assume $d>0$ and $k\le d$. Let
    $T:\MN{d}\to\MN{d'}$ have Choi matrix $\tau$. Suppose that
    $R_P\tau R_P^\dagger\geq0$ for every Hermitian projection
    $P\in\MN{d}$ of rank $k$. Then $T$ is $k$-positive.
\end{theorem}

\begin{proof}
    \leanok
    By Theorem~\ref{thm:k_positive_iff_right_choi_compressions}, it is enough
    to prove positivity of the right-factor compression by an arbitrary
    $X\in M_{d\times k}(\C)$. Choose a rank-$k$ Hermitian projection $P$ with
    $PX=X$ using
    Lemma~\ref{lem:rank_k_projection_fixing_rectangular_factor}. The assumed
    positivity of $R_P\tau R_P^\dagger$, together with
    Lemma~\ref{lem:square_sandwich_fixed_rectangular_compression}, gives
    positivity of the rectangular compression by $X$.
\end{proof}

\begin{theorem}[Rank-$k$ input-projection form of the rectangular Choi criterion]
    \label{thm:rank_k_right_projection_choi_criterion}
    \lean{ChoiJamiolkowski.isNPositiveMap_iff_forall_rankProjection_choiMatrix_sandwich_posSemidef_rectangular}
    \leanok
    \uses{thm:rank_k_right_projection_choi_sandwich_forward,
        thm:rank_k_right_projection_choi_sandwich_converse}
    Assume $d>0$ and $k\le d$. A linear map
    $T:\MN{d}\to\MN{d'}$ is $k$-positive if and only if, for every
    Hermitian projection $P\in\MN{d}$ of rank $k$,
    $R_P\tau R_P^\dagger\geq0$, where $\tau$ is the Choi matrix of $T$.
    This is the rank-$k$ projection formulation of
    \cite[Chapter~3, Proposition~3.1]{Wolf2012Quantum}.
\end{theorem}

\begin{proof}
    \leanok
    The forward direction is
    Theorem~\ref{thm:rank_k_right_projection_choi_sandwich_forward}. The
    reverse direction is
    Theorem~\ref{thm:rank_k_right_projection_choi_sandwich_converse}.
\end{proof}

\section{Trace adjoints and the positivity hierarchy}

The trace-pairing adjoint was introduced in
Definition~\ref{def:trace_pairing_adjoint}. Its compatibility with
ampliation makes $k$-positivity invariant under trace adjoints and identifies
the endpoints of the positivity hierarchy. Routine closure properties of the
cones of $k$-positive maps are collected in
Section~\ref{sec:app_positive_not_cp_k_positive_closure}.

\begin{theorem}[Trace-pairing adjoint of an ampliation]
    \label{thm:trace_pairing_adjoint_ampliation}
    \lean{nPositiveAmpliation_traceAdjointMap}
    \leanok
    \uses{def:blockwise_ampliation, def:trace_pairing_adjoint,
        thm:trace_pairing_adjoint_identity}
    For every linear map $E : \MN{D}\to\MN{D}$, the trace-pairing adjoint
    commutes with ampliation:
    \begin{align}
        (E^{(k)})^*
         & =(E^*)^{(k)}.
        \label{eq:schwarz_trace_adjoint_ampliation}
    \end{align}
\end{theorem}

\begin{proof}
    \leanok
    For block matrices write $\rho_{pq}$ and $X_{pq}$ for their
    $\MN{D}$-valued blocks. Expanding the trace by blocks and applying
    (\ref{eq:schwarz_trace_adjoint_pairing}) gives
    \begin{align}
        \tr((E^*)^{(k)}(\rho)X)
         & =\sum_{p,q}\tr(E^*(\rho_{pq})X_{qp})
        =\sum_{p,q}\tr(\rho_{pq}E(X_{qp}))
        =\tr(\rho E^{(k)}(X)).
        \notag
    \end{align}
    Nondegeneracy of the trace pairing identifies the two adjoints, proving
    (\ref{eq:schwarz_trace_adjoint_ampliation}).
\end{proof}

\begin{theorem}[Trace-pairing adjoint of a $k$-positive map]
    \label{thm:k_positive_trace_pairing_adjoint}
    \lean{IsNPositiveMap.traceAdjointMap}
    \leanok
    \uses{def:k_positive_map, thm:trace_pairing_adjoint_ampliation,
        thm:trace_pairing_adjoint_positive}
    If $E$ is $k$-positive, then $E^*$ is $k$-positive.
\end{theorem}

\begin{proof}
    \leanok
    Since $E$ is $k$-positive, its ampliation $E^{(k)}$ is positive. The
    trace-pairing adjoint of this positive map is positive, and
    (\ref{eq:schwarz_trace_adjoint_ampliation}) identifies that adjoint with
    $(E^*)^{(k)}$.
\end{proof}

\begin{theorem}[$k$-positivity is trace-adjoint invariant]
    \label{thm:k_positive_trace_pairing_adjoint_iff}
    \lean{isNPositiveMap_traceAdjointMap_iff}
    \leanok
    \uses{thm:k_positive_trace_pairing_adjoint,
        thm:trace_pairing_adjoint_involutive}
    A linear map is $k$-positive if and only if its trace-pairing adjoint is
    $k$-positive.
\end{theorem}

\begin{proof}
    \leanok
    One direction is
    Theorem~\ref{thm:k_positive_trace_pairing_adjoint}. The other follows by
    applying the same theorem to $E^*$ and using
    Theorem~\ref{thm:trace_pairing_adjoint_involutive}.
\end{proof}

\begin{theorem}[Completely positive maps are $k$-positive]
    \label{thm:cp_implies_k_positive}
    \lean{IsCPMap.isNPositiveMap}
    \leanok
    \uses{def:cp_map}
    Every completely positive map is $k$-positive for all $k \ge 1$.
\end{theorem}

\begin{proof}
    \leanok
    Choose a Kraus representation $E(X)=\sum_i K_iXK_i^\dagger$. For every
    $Y\succeq0$ in $\MN{D}\otimes\MN{k}$,
    \begin{align}
        (E\otimes\id_k)(Y)
         & =\sum_i (K_i\otimes\Id_k)Y(K_i\otimes\Id_k)^\dagger
        \succeq0.
        \notag
    \end{align}
    Thus every ampliation $E\otimes\id_k$ is positive.
\end{proof}

\begin{theorem}[One-positive maps are positive]
    \label{thm:one_positive_iff_positive}
    \lean{isNPositiveMap_one_iff_isPositiveMap}
    \leanok
    \uses{def:k_positive_map, def:positive_map}
    A linear map is $1$-positive if and only if it is positive.
\end{theorem}

\begin{proof}
    \leanok
    The ampliation by $\MN{1}$ is identified with the original matrix algebra:
    all block indices have the unique value, so positivity of
    $E\otimes\id_1$ is exactly positivity of $E$.
\end{proof}

\begin{theorem}[$D$-positive maps on nonzero $\MN{D}$ are completely positive]
    \label{thm:d_positive_iff_completely_positive}
    \lean{isCPMap_iff_isNPositiveMap_card}
    \leanok
    \uses{def:k_positive_map, def:cp_map, thm:cp_iff_choi_posSemidef}
    For maps $E:\MN{D}\to\MN{D}$ with $D\ge 1$, complete positivity is
    equivalent to $D$-positivity.
\end{theorem}

\begin{proof}
    \leanok
    Complete positivity implies $k$-positivity for every $k$. Conversely,
    assume $D\ge 1$. If $E$ is $D$-positive, apply $E\otimes\id_D$ to the
    maximally entangled projector $\ketbra{\Omega}{\Omega}$. The result is
    the Choi matrix of $E$, hence it is positive semidefinite. Choi's theorem
    then gives complete positivity.
\end{proof}
